% deixis — the whitepaper. Theory backing docs/TREE.md, SLOTS.md, PATH.md and the
% design records 0001–0005. Plain LaTeX: article + AMS packages only, so any
% distribution (TeX Live, MiKTeX, tectonic) compiles it without ceremony.
\documentclass[11pt]{article}

\usepackage[margin=1.15in]{geometry}
\usepackage{amsmath,amssymb,amsthm}
\usepackage{booktabs}
\usepackage[hidelinks]{hyperref}
\hypersetup{pdfauthor={Julian Matschinske (julian@matschinske.com)},
  pdftitle={deixis: A Minimal Floor for Keyed Tree Structure}}
\usepackage{enumitem}

\newtheorem{theorem}{Theorem}[section]
\newtheorem{lemma}[theorem]{Lemma}
\newtheorem{proposition}[theorem]{Proposition}
\newtheorem{corollary}[theorem]{Corollary}
\theoremstyle{definition}
\newtheorem{definition}[theorem]{Definition}
\newtheorem{example}[theorem]{Example}
\newtheorem{remark}[theorem]{Remark}
\newtheorem{stipulation}[theorem]{Stipulation}

% ----- notation ---------------------------------------------------------------
\newcommand{\Bytes}{\mathbb{B}}                    % finite byte strings
\newcommand{\NodeT}[1]{\mathrm{Node}(#1)}
\newcommand{\Leaf}{\mathrm{Leaf}}
\newcommand{\Struct}{\mathrm{Struct}}
\newcommand{\dom}{\mathrm{dom}\,}
\newcommand{\eqS}{=_{\approx}}                     % lifted identity
\newcommand{\finmap}[2]{#1 \rightharpoonup_{\mathrm{fin}} #2}
\newcommand{\lex}{<_{\mathrm{lex}}}
\newcommand{\magn}{\mathrm{mag}}
\newcommand{\enc}{\mathsf{enc}}
\newcommand{\encD}{\mathsf{enc}_{\!D}}
\newcommand{\Paths}{\mathrm{Paths}}
\newcommand{\VO}{V_{\mathrm{o}}}                   % ontos values
\newcommand{\Atom}{\mathrm{Atom}}
\newcommand{\Tup}{\mathrm{Tuple}}
\newcommand{\eqO}{=_{\mathrm{o}}}

\title{deixis: A Minimal Floor for Keyed Tree Structure\\
  \large Identity, pointing, and readings over one container}
\author{Julian Matschinske \quad
  \href{mailto:julian@matschinske.com}{\small\texttt{julian@matschinske.com}}}
\date{September 2026 — draft v4.1, mandatory node values (ADR 0010)}

\begin{document}
\maketitle

\begin{abstract}
deixis is a specification of finite byte-keyed trees over one opaque slot:
$\NodeT{T} = T \times (\Bytes \rightharpoonup_{\mathrm{fin}} \NodeT{T})$, so that every
node carries exactly one own value and a finite map of children; where absence is wanted,
it is a choice of slot. It fixes
the structural half of a data floor --- shape, keys, paths, canonical order --- and
commits to nothing about what a tree holds. Construction and navigation need only a
carrier; identity additionally uses a supplied equivalence. This paper gives the theory the specification
rests on. We define the tree as the initial algebra of a container functor, lift a
slot equality to a node identity and prove the lifting a congruence, characterize
identity by pointing (paths and the values at them), and give resolution as a partial
action of the free monoid of keys. We prove that a node is exactly its parts, and that
every node can gain a child under a fresh key while every existing path and value is
kept. The model of this paper's earlier drafts, in which a node either holds a value or
has children, embeds into the optional instantiation $\NodeT{\mathrm{Option}(T)}$, and
for any inhabited slot no map in the other direction preserves paths and own values. We isolate the admission conditions a slot must meet per capability
tier, proving that a byte-valued complete invariant exists iff the slot quotient is
countable, and that existence and computability come apart. We prove the three laws of the positional
spelling $\kappa$ (injective, prefix-free, order-preserving), from which sequence
semantics embed into the key space at zero interpretive cost. We then develop the
reading principle: collection and positional semantics (order, sets, multiplicity,
records) enter as \emph{structure profiles} --- embedding/retraction pairs into the
key space --- rather than as constructors or type parameters, and we prove the
transport laws for the sequence and set profiles at their stated scope. Finally we
exhibit the ontos value model as an exact profile of
$\NodeT{\mathrm{Option}(\Bytes)}$. Its mathematical projection and recognizer extend
the earlier bridge; the shipping bridge still targets the previous core API.
\end{abstract}

\medskip
\noindent\fbox{\parbox{\dimexpr\linewidth-2\fboxsep-2\fboxrule\relax}{\small\raggedright
\textbf{Scope of draft v4 (2026-09-23).} This paper states ADR~0010's mandatory-value
floor, $\NodeT{T} = T \times (\finmap{\Bytes}{\NodeT{T}})$. The previous optional floor
is its explicit instantiation at $\mathrm{Option}(T)$; the earlier leaf/structure
model embeds into that instantiation (Remark~\ref{rem:previous}). The proofs are
mathematical arguments, not proof-assistant verification. The encoder in
\S\ref{sec:codec} is an abstract construction demonstrating the lifted invariant
law, not a specification of published bytes. The candidate \texttt{CODEC.md} now
states \texttt{deixis-codec-v2}, a grammar for this model (ADR~0011); its proofs,
vectors and freeze obligations are separate from this paper. Consumer implementations
are not certified by these structural proofs.}}

\clearpage
\section{Introduction}

\paragraph{Mathematical scope and review.}
The aim is expressive power through a small, coherent set of concepts. Defined
objects have their defining properties; constructions and derived laws are
assessed within their stated domains and assumptions. There is no failure until
the model defines a failure concept. Partial selection and expressly defined
recognizer refusals remain part of the mathematics; physical contingencies enter
only through an explicit model that introduces them. A calculator losing power
does not qualify $2+2=4$, and an unmodeled stuck tape does not qualify a Turing
machine's transition rule.

Review this paper for internal consistency, precise definitions, well-defined
constructions, valid derivations, expressive reach, and economy of concepts.
A defining law states what an admitted object is; a claim that a particular
construction supplies such an object still needs justification. Counterexamples
must satisfy the stated premises. Do not add feedback about phenomena the model
does not contain, or make implementation a prerequisite for theoretical value.
Constructions, embeddings, and composition results can establish that value
mathematically. Separate implementation claims concern whether particular
artifacts realize the model. Draft v4.1 adds this methodological clarification to
v4 without changing its definitions or theorems; the
\href{https://github.com/Bitspark/deixis/blob/main/docs/paper/README.md}{paper's README}
states the common scope.

Systems that let independent parties agree without trusting each other need, at the
bottom, a notion of \emph{value} whose identity can be decided by looking --- no
context, no clock, no authority. Every layer above such a floor (canonical encodings,
content addressing, re-checkable derivations, typed contracts, revocable mandates)
inherits its identity from below and can weaken but never repair it. The floor is
therefore where stipulation is cheapest and where every avoidable commitment is
dearest: a constructor added to the floor is inherited by everyone forever; a meaning
attached to the floor is an opinion everyone must share.

deixis is a floor designed under that discipline, one level beneath a byte-atom value
model in production use (ontos, \S\ref{sec:ontos}). Its design bets are:

\begin{enumerate}[leftmargin=1.6em]
\item \textbf{One container, closed.} A node is an own value and a finite map from byte
  keys to nodes. There is one constructor and no node kind, no edge metadata
  beyond the key itself, no ordered sibling list --- and the collection and positional readings a consumer
  reaches for are recovered as proved embeddings (\S\ref{sec:profiles}; the general
  reading principle is stated there with its current scope, not as a universality
  theorem).
\item \textbf{One hole, minimal interface.} The slot, the own value every node carries,
  admits any carrier whatsoever, an optional one included; deixis interacts with it through a single supplied relation, and each
  further capability (decidability, canonical bytes, decodability) is priced as one
  further supplied operation (\S\ref{sec:slots}).
\item \textbf{Keys are bytes, not a parameter.} Byte keys discharge, once and for
  everyone, the two obligations any key sort owes: an exact decidable equality (for
  identity) and a total order (for canonicalization). This is an engineering economy
  rather than a mathematical necessity (Remark~\ref{rem:keys}); what keys \emph{mean}
  is left to readings above the floor (\S\ref{sec:pos}).
\item \textbf{Structural identity is fixed; its value clause is supplied.} deixis
  fixes exact key domains and pointwise recursion, and parameterizes exactly the
  comparison of own values by the supplied equivalence. That no layer above
  may coarsen or refine the result is a \emph{protocol rule} the ecosystem enforces on
  its observers (Remark~\ref{rem:extensional}), not a property of the raw type.
\end{enumerate}

The contributions of this paper are the proofs: that the lifting is a congruence
(Thm.~\ref{thm:equiv}, \ref{thm:congruence}); that identity is exactly co-pointing
(Thm.~\ref{thm:pointing}); that resolution is a partial monoid action respecting
identity (Prop.~\ref{prop:action}); that a node is exactly its parts and that split
and plug are inverse (Prop.~\ref{prop:parts}, \ref{prop:splitplug}); that fresh
attachment is exact and available at every node (Thm.~\ref{thm:attach},
Cor.~\ref{cor:attach}), and that for an inhabited slot no path- and value-preserving map
back to the previous model exists (Prop.~\ref{prop:nonembed}); the countability boundary for complete invariants
and its computability gap (Thm.~\ref{thm:countable}, Prop.~\ref{prop:gap}); the
correctness of the canonical encoding combinator (Thm.~\ref{thm:codec}); the three
$\kappa$ laws (Thm.~\ref{thm:kappa}); the transport laws for the sequence and set
profiles (Prop.~\ref{prop:pos-transport}, \ref{prop:set-transport}); and the exactness
of the ontos projection (Thm.~\ref{thm:ontos}).

\section{The tree}\label{sec:tree}

\emph{Metatheory.} The semantic results of this paper (\S\ref{sec:tree}--%
\S\ref{sec:profiles}, \S\ref{sec:ontos}) are stated in ordinary classical set theory:
sets, quotients, and setoids as carriers-with-equivalence, with no constructive or
computability commitments. Statements about \emph{effectiveness} are separate and
explicit: they are read in ordinary computability over $\mathbb{N}$ and $\Bytes$,
algorithms consume the named representations (byte strings; program indices where
indicated), and partiality conventions are stated where used
(Prop.~\ref{prop:gap} spells out its two readings of ``computable on $T$''). No
represented-space machinery is invoked, and no claim in this paper depends on one.

Fix $\Bytes$, the set of finite octet strings, with decidable equality and the total
unsigned-lexicographic order $\lex$. For sets $K, X$ write $\finmap{K}{X}$ for the set
of finite partial maps $K \rightharpoonup X$.

Write $\mathrm{Option}(T) = \mathbf{1} \sqcup T$, with elements $\mathrm{None}$ and
$\mathrm{Some}(x)$ for $x \in T$. It is not part of the floor. It is one choice of slot
(Remark~\ref{rem:instances}), used below for the previous model and for the container
profiles.

\begin{definition}[The functor]\label{def:functor}
For a set $T$, let $F_T(X) \;=\; T \times (\finmap{\Bytes}{X})$.
\end{definition}

Since $\finmap{\Bytes}{X} \cong \sum_{K \in \mathcal{P}_{\mathrm{fin}}(\Bytes)} X^{K}$,
$F_T(X) \cong \sum_{(o, K)} X^{K}$ over $(o, K) \in T \times
\mathcal{P}_{\mathrm{fin}}(\Bytes)$: $F_T$ is a container (polynomial) functor with
shapes $T \times \mathcal{P}_{\mathrm{fin}}(\Bytes)$ and positions $K$
at shape $(o, K)$. Its initial algebra exists and is the corresponding W-type
\cite{martinlof,aag}.

\begin{definition}[Nodes]\label{def:node}
$\NodeT{T}$ is the initial $F_T$-algebra. We write its single constructor
$\mathrm{Node} : T \times (\finmap{\Bytes}{\NodeT{T}}) \to \NodeT{T}$.
In $\mathrm{Node}(o, m)$, $o \in T$ is the node's \emph{own value} and $m$ its
\emph{children}.
\end{definition}

Every node carries exactly one own value and a finite map of children, which may be
empty. A node with children keeps its own value, and a node without children still has
one. There are no node kinds: ``leaf'' names a node without children, which is a fact
about shape and says nothing about the value. A missing node is a result of navigation
(\S\ref{sec:path}), never a special value of $T$.

Initiality gives the two structural facts the specification stipulates rather than
derives, and it is worth recording why stipulation is nevertheless the honest
presentation: the specification addresses implementations in languages whose type
systems do not exhibit initial algebras, so the floor asserts the consequences
directly and this paper certifies that a model exists in which all of them hold.

\begin{stipulation}[Structural facts]\label{stip:structural}
(i) \emph{Injectivity}: $\mathrm{Node}$ is injective.
(ii) \emph{Finiteness and well-foundedness}: every node is a finite tree; there is no
infinite descent through children, and structural induction and recursion are valid.
\end{stipulation}

Both hold in the initial algebra. The initial algebra's structure map is injective
by Lambek's lemma, and W-types are well-founded by construction. Note what is \emph{not}
present: no node kinds, no node metadata beyond the own value, no edge labels, and no
ordering among siblings. The only degrees of freedom of a node are its own value, its key
set, and its children.

\begin{remark}[Optionality is a slot choice]\label{rem:instances}
The floor imposes no absence convention. A consumer that wants one instantiates the slot
with it. $\NodeT{\mathrm{Option}(A)}$ carries an explicit $\mathrm{None}$ or
$\mathrm{Some}(a)$ at every node, and under the usual tag-respecting equality
$\mathrm{None} \neq \mathrm{Some}(a)$ even when $a$ is ``empty''. That inequality belongs
to the slot and is declared by its supplier; the floor does not stipulate it.
$\NodeT{\mathbf{1}}$ is pure keyed shape. $\NodeT{\mathbf{0}}$ is empty, since every tree
has a root value. In every instantiation a sentinel inside the slot is a value at an
existing node, never the absence of a node.
\end{remark}

\begin{remark}[The previous model]\label{rem:previous}
Earlier drafts of this paper took $F_T(X) = T \sqcup (\finmap{\Bytes}{X})$, with
constructors $\Leaf$ and $\Struct$; write $L(T)$ for its initial algebra. It embeds into
this model at the slot $\mathrm{Option}(T)$, not at $T$ itself, by
$E : L(T) \to \NodeT{\mathrm{Option}(T)}$ with
\[
E(\Leaf(x)) = \mathrm{Node}(\mathrm{Some}(x), \emptyset), \qquad
E(\Struct(m)) = \mathrm{Node}(\mathrm{None},\, k \mapsto E(m(k))).
\]
The image of $E$ is exactly the trees over $\mathrm{Option}(T)$ in which $\mathrm{Some}$
occurs only at nodes without children. That was the previous model's one restriction: a
node could not hold a value and have children. No embedding is promised at an unchanged
slot $T$. An intermediate design (ADR~0009 of the specification) made
$\mathrm{Option}(T) \times (\finmap{\Bytes}{X})$ the floor itself; that model is exactly
$\NodeT{\mathrm{Option}(T)}$, without the restriction. In the previous model a node holding
a value could gain a child only by being replaced, which moved its value to another path.
\end{remark}

\begin{remark}[Keys are not a parameter]\label{rem:keys}
A two-parameter floor $\mathrm{Node}[K, T]$ was considered and rejected in the
specification's first design record. The honest form of the argument is an obligation
count, not an impossibility: a key sort $K$ could, like the slot, owe only an equality
at the identity tier and a lawful encoder at the canonicalization tier --- the
encoder's bytes would then \emph{supply} the total order. What fixing $K = \Bytes$
buys is that both obligations are discharged once, by the type itself, for every
consumer forever, and the one-hole notation survives. An engineering economy,
deliberately taken; \S\ref{sec:pos} shows nothing positional is lost.

The total order also does quiet foundational work. With a decidable total order,
$\mathcal{P}_{\mathrm{fin}}(\Bytes)$ has canonical representatives --- strictly
sorted, duplicate-free entry lists --- so the branching functor is a \emph{strict}
polynomial functor rather than a quotient of one: $\NodeT{T}$ mechanizes as a plain
inductive type carrying a decidable, hence proof-irrelevant, sortedness invariant,
and no \emph{structural} setoid arises before the slot setoid is even in play. (This
is exactly the shipped representation: sorted-unique entry vectors.) For an
arbitrary key sort with mere decidable equality, finite maps are a quotient of entry
lists up to reordering, and every mechanization would inherit setoid structure at
the structural level.
\end{remark}

\section{Identity}\label{sec:identity}

deixis does not define identity; it lifts one. The slot is not a type but a
\emph{setoid}: a pair $(T, \approx)$ of a carrier and a relation. Nothing else about
the slot is ever consulted.

\begin{definition}[Lifted identity]\label{def:lift}
Given $(T, \approx)$, define $\eqS$ on $\NodeT{T}$ by structural recursion:
\[
\mathrm{Node}(o, m) \eqS \mathrm{Node}(o', n) \iff
  o \approx o' \;\wedge\; \dom m = \dom n \;\wedge\; \forall k \in \dom m.\;
  m(k) \eqS n(k).
\]
Key comparison in $\dom m = \dom n$ is octet equality; the recursion is well-founded
by Stipulation~\ref{stip:structural}(ii).
\end{definition}

\begin{theorem}[The lifting preserves equivalence]\label{thm:equiv}
If $\approx$ is reflexive (resp.\ symmetric, transitive), so is $\eqS$. In particular
if $(T,\approx)$ is a setoid, $(\NodeT{T}, \eqS)$ is a setoid.
\end{theorem}

\begin{proof}
Each property of $\eqS$ separately, by induction on the node, which is one case since
there is one constructor. Reflexivity: $\mathrm{Node}(o, m) \eqS \mathrm{Node}(o, m)$
needs $o \approx o$, which is reflexivity of $\approx$,
$\dom m = \dom m$, and $m(k) \eqS m(k)$, the last by the inductive hypothesis on the
strictly smaller $m(k)$. Symmetry and transitivity are the same induction, using the
same property of $\approx$ on own values, symmetry and transitivity of set equality on
domains, and
the inductive hypothesis pointwise. (The converse direction of the slogan ``deixis lifts
what it is given and does not repair it'' is immediate: if $x \not\approx x$ then
$\mathrm{Node}(x, \emptyset) \not\eqS \mathrm{Node}(x, \emptyset)$.)
\end{proof}

\begin{remark}[The admission rule: three classes of maps]\label{rem:extensional}
The raw type $\NodeT{T}$ still distinguishes what $\eqS$ merges:
$\mathrm{Node}(x, \emptyset) \neq \mathrm{Node}(y, \emptyset)$
as elements even when $x \approx y$, and an observer that inspects raw slot
representations can tell them apart. Every ``merged is merged forever'' statement in
this paper is therefore relative to \emph{extensional} observers --- functions that
are setoid morphisms, respecting $\approx$ and $\eqS$.

The specification's rule that no layer may coarsen or refine identity has two halves
with \emph{different} enforcers, and an earlier draft of this remark conflated them.
The no-\emph{refinement} half is exactly extensionality: a setoid morphism cannot
separate what $\eqS$ merges. The no-\emph{coarsening} half is \textbf{not} ---
a constant observer is extensional and maximally coarsening. Coarsening is instead
excluded, where it must be, by requiring the map to \emph{reflect} equality as well
as preserve it: $f(N) = f(M) \Rightarrow N \eqS M$. The ecosystem accordingly admits
three classes of maps, with graded obligations:

\begin{enumerate}[label=(\alph*), leftmargin=2.2em]
\item \emph{Observations} --- extensional only. A size, a leaf count, a predicate is
  naturally lossy; nothing downstream treats its output as an identity surrogate.
\item \emph{Identity-bearing representations and migrations} --- extensional
  \emph{and} equality-reflecting; equivalently, embeddings of quotient setoids. The
  lawful-invariant law \eqref{eq:codeclaw}, profile identity transport
  (Def.~\ref{def:profile}(iii)), and the ontos embedding
  (Thm.~\ref{thm:ontos}(iii)) are exactly this two-sided condition. Reflection does
  not survive post-composition: a pipeline's membership in this class must be
  re-established end to end, or composed from certified reflecting stages --- the
  label is a claim about a whole map, never an inheritance from its parts.
\item \emph{Recognizers and decoders} --- exact partial inverses.
  Definition~\ref{def:profile}(i),(ii) states their laws; they refuse rather than repair.
\end{enumerate}

Raw observation is not a fourth class but the \emph{boundary} of the taxonomy, and the
specification permits it. An own-value accessor may inspect the representative;
that is what makes a debugger, a wire dump or a diagnostic possible at all, and a doctrine
forbidding it would forbid tooling. What such an observer may not do is launder its result.
A quantity that depends on the representative is not a function of the floor value, and the
obligation attaches to the \emph{claim} rather than to the inspection: a
representative-sensitive result may never be presented as value-determined. Where a map is
offered as a function of the value it must be $\approx$-congruent, which places it in class
(a); where it is not, it must say which it is. This is the discipline of (c) stated for
observers rather than for decoders --- the failure mode is not looking, it is looking and
then claiming more than was seen.

This is a protocol obligation the ecosystem enforces (and its conformance vectors
test), not a property the raw data type can enforce by itself. Every operation
defined in this paper is admitted under its class: the observations of
\S\ref{sec:tree}--\ref{sec:path} are extensional (Prop.~\ref{prop:action}(ii) is the
case proved for resolution), and the maps this paper treats as identity-bearing are
proved to reflect.
\end{remark}

\begin{definition}[Replacement]\label{def:replace}
For a path $p \in \Bytes^{*}$ (Def.~\ref{def:path}) with $N/p$ defined, let
$N[p \mapsto M]$ be the node obtained by replacing the subtree at $p$ with $M$,
defined by recursion on $p$: $N[\varepsilon \mapsto M] = M$, and
$\mathrm{Node}(o, m)[(k \cdot p) \mapsto M] = \mathrm{Node}(o, m[k \mapsto m(k)[p \mapsto
M]])$. Every node on the way down keeps its own value $o$ and its other children; only
the subtree at $p$ is replaced, own value included.
\end{definition}

\begin{theorem}[Congruence]\label{thm:congruence}
Let $\approx$ be an equivalence. If $N/p$ is defined and $M \eqS M'$, then
$N[p \mapsto M] \eqS N[p \mapsto M']$. More generally $\eqS$ is a congruence for
$\mathrm{Node}$: if $o \approx o'$, $\dom m = \dom n$ and $m(k) \eqS n(k)$ pointwise
then $\mathrm{Node}(o, m) \eqS \mathrm{Node}(o', n)$.
\end{theorem}

\begin{proof}
The general statement is Definition~\ref{def:lift} read as an introduction rule. The
replacement statement follows by induction on $p$: for $\varepsilon$ it is the
hypothesis; for $k \cdot p$, the two results share the own value $o$ of the node being
descended through, which is related to itself by reflexivity of $\approx$, and their child
maps agree on every key except $k$, where the inductive hypothesis applies. Reflexivity
(Thm.~\ref{thm:equiv}) covers the untouched entries.
\end{proof}

\subsection{Functoriality in the slot}

The slot is a parameter, and the parametricity is itself lawful: $\NodeT{-}$ is a
functor, and the identity theory transports along it.

\begin{definition}[Map]\label{def:map}
For $f : T \to U$ define $\mathrm{map}\,f : \NodeT{T} \to \NodeT{U}$ by structural
recursion:
\[
\mathrm{map}\,f\,(\mathrm{Node}(o, m)) = \mathrm{Node}(f\,o,\; k \mapsto \mathrm{map}\,f\,(m(k))).
\]
Keys are untouched. $\mathrm{map}$ applies $f$ to the own value of \emph{every} node, and
never adds or removes a node. At $T = \mathrm{Option}(A)$ it may therefore change
$\mathrm{None}$; the presence-preserving mapping is the explicit specialization
$\mathrm{map}\,(\mathrm{Option}.\mathrm{map}\,g)$.
\end{definition}

\begin{proposition}[Functor laws; structure invariance; setoid transport]\label{prop:functor}
\emph{(i)} $\mathrm{map}\,\mathrm{id} = \mathrm{id}$ and
$\mathrm{map}\,(g \circ f) = \mathrm{map}\,g \circ \mathrm{map}\,f$.
\emph{(ii)} $\Paths(\mathrm{map}\,f\,N) = \Paths(N)$, and resolution commutes:
$(\mathrm{map}\,f\,N)/p \simeq \mathrm{map}\,f\,(N/p)$ (Kleene equality).
\emph{(iii)} If $f$ is a setoid morphism --- $x \approx_T y \Rightarrow
f(x) \approx_U f(y)$ --- then $N =_{\approx_T} M \Rightarrow
\mathrm{map}\,f\,N =_{\approx_U} \mathrm{map}\,f\,M$. Hence
$(T, \approx) \mapsto (\NodeT{T}, \eqS)$ with $f \mapsto \mathrm{map}\,f$ is an
endofunctor on the category of setoids, of which Theorem~\ref{thm:equiv} is the
object part.
\end{proposition}

\begin{proof}
All three by structural induction, which has one case. (i) The own values agree because
$\mathrm{id}\,o = o$ and $(g \circ f)\,o = g\,(f\,o)$, and the children agree pointwise
by the inductive hypothesis. (ii) $\mathrm{map}$ rewrites no key, so path sets coincide;
the resolution clause follows by induction on $p$.
(iii) Own values: related by the morphism property. Children: untouched, equal domains and pointwise-related subtrees by the
inductive hypothesis --- exactly Definition~\ref{def:lift}.
\end{proof}

\begin{corollary}[Embeddings lift]\label{cor:embed}
Let $f$ be a \emph{setoid embedding}: it preserves equality ($x \approx_T y
\Rightarrow f(x) \approx_U f(y)$) \emph{and} reflects it ($f(x) \approx_U f(y)
\Rightarrow x \approx_T y$) --- the identity-bearing class of
Remark~\ref{rem:extensional}. Then $\mathrm{map}\,f$ is a setoid embedding as well:
preservation is Prop.~\ref{prop:functor}(iii), and for reflection, by the pointing
characterization (Thm.~\ref{thm:pointing}),
$\mathrm{map}\,f\,N =_{\approx_U} \mathrm{map}\,f\,M$ forces equal path sets (untouched by
$\mathrm{map}$) and $\approx_U$-equal images at every path, which reflection returns to
$\approx_T$-equal own values. Slot-level embeddings
lift to node-level embeddings; neither direction alone would suffice for the claim.
\end{corollary}

\begin{theorem}[Lifting commutes with quotienting]\label{thm:quotient}
Let $q : T \to T/\!\approx$ be the projection onto the quotient setoid (whose
equality is discrete equality of classes). Then $\mathrm{map}\,q$ induces an
isomorphism of setoids
\[
\NodeT{T}/\!\eqS \;\cong\; \NodeT{T/\!\approx}.
\]
\end{theorem}

\begin{proof}
$q$ preserves and reflects equality by definition of the quotient, so
$\mathrm{map}\,q$ is a setoid embedding (Cor.~\ref{cor:embed}) and descends to an
injection $\NodeT{T}/\!\eqS \hookrightarrow \NodeT{T/\!\approx}$. Surjectivity:
given $M \in \NodeT{T/\!\approx}$, choose a representative for each of its finitely
many value classes, one per path --- finite choice, a theorem of ZF, within the classical
metatheory of \S\ref{sec:tree}; at the decodable tier no choice is needed, the
exact partial decoder supplying the representative (cf.\
Prop.~\ref{prop:set-transport}) --- and the resulting $N \in \NodeT{T}$ has
$\mathrm{map}\,q\,N = M$ by structural induction.
\end{proof}

\begin{remark}[The protocol stance is the quotient semantics]\label{rem:hofmann}
Theorem~\ref{thm:quotient} is the precise sense in which the admission discipline of
Remark~\ref{rem:extensional} \emph{is} quotient semantics rather than a substitute
for it: the setoid model interprets quotient types \cite{hofmann}, and every
extensional observer of $\NodeT{T}$ factors uniquely through $\NodeT{T/\!\approx}$.
A mechanization in a prover with quotient types can work on the quotient side and
transport every result across this isomorphism.
\end{remark}

\begin{remark}[Transport of profiles]
A structure profile (\S\ref{sec:profiles}) is a condition on keys, so
$\mathrm{map}\,f$ preserves a positional image verbatim. A profile whose recognizer
consults own values --- the set profile's coherence $k = \enc(x)$ for a member
$\mathrm{Node}(\mathrm{Some}(x), \emptyset)$ of its optional carrier --- transports
along $\mathrm{map}\,(\mathrm{Option}.\mathrm{map}\,f)$ exactly when the member encoders
agree, $\enc_U \circ f = \enc_T$;
otherwise the image is honestly left rather than repaired.
\end{remark}

\subsection{Identity is pointing}

\begin{definition}[Paths of a node]\label{def:paths}
$\Paths(N) \subseteq \Bytes^{*}$:\; $\varepsilon \in \Paths(N)$ always, and
$k \cdot p \in \Paths(\mathrm{Node}(o, m))$ iff $k \in \dom m$ and $p \in \Paths(m(k))$.
Write $\mathrm{val}_N(p) \in T$ for the own value of $N/p$; it is defined at every
$p \in \Paths(N)$. A path is \emph{terminal} iff $N/p$ has no children, and terminal and
interior paths alike carry a value.
\end{definition}

\begin{theorem}[Pointing characterization]\label{thm:pointing}
Let $\approx$ be an equivalence. $N \eqS M$ iff
\emph{(a)} $\Paths(N) = \Paths(M)$, and
\emph{(b)} $\mathrm{val}_N(p) \approx \mathrm{val}_M(p)$ for every $p \in \Paths(N)$.
\end{theorem}

\begin{proof}
Write $N = \mathrm{Node}(o, m)$ and $M = \mathrm{Node}(o', n)$.
($\Rightarrow$) By induction on $N$. At $\varepsilon$, $o \approx o'$ is (b) at
$\varepsilon$. Then $\dom m = \dom n$ makes the length-one prefixes agree, and the
inductive hypothesis on each $m(k) \eqS n(k)$ extends (a) and (b) under the prefix $k$.
($\Leftarrow$) By induction on $N$. (b) at $\varepsilon$ is $o \approx o'$. Next, $k \in \dom m$ iff $k \in \Paths(N)$ iff (by (a)) $k \in \Paths(M)$
iff $k \in \dom n$, so the domains agree. For each shared $k$, the sets
$\{p : k \cdot p \in \Paths(\cdot)\}$ and the restricted (b) are exactly conditions
(a) and (b) for $m(k), n(k)$, and the inductive hypothesis yields $m(k) \eqS n(k)$.
\end{proof}

\begin{example}[Neither condition is redundant]
For any $x \in T$, the nodes
\[
\mathrm{Node}(x, \{a \mapsto \mathrm{Node}(x, \emptyset)\})
\quad\text{and}\quad \mathrm{Node}(x, \emptyset)
\]
carry related values wherever both have a path, but their path sets differ ($a$ is in
one and not the other). Conversely, for $x \not\approx y$, $\mathrm{Node}(x, \emptyset)$
and $\mathrm{Node}(y, \emptyset)$ have the same path set $\{\varepsilon\}$ and are
separated only by (b). At $T = \mathrm{Option}(A)$ the first pair can carry
$\mathrm{None}$ throughout: a path whose value reads as absence is still a path.
\end{example}

The theorem is the origin of the name: a node is what it points at, and nothing else.

\begin{remark}[What identity is not]
Two non-theorems are design content. $\eqS$ is not \emph{object} identity: no address,
history, or sharing enters (sharing of equal subtrees is representation and
unobservable). And $\eqS$ is not coarser than the lifting: no equation such as
``a one-child node equals its child'' or ``a node whose value is some empty sentinel
equals no node at all'' is imposed; each would import a consumer's
theory into the floor.
\end{remark}

\section{Resolution}\label{sec:path}

\begin{definition}[Paths and resolution]\label{def:path}
$\mathrm{Path} = \Bytes^{*}$, the free monoid under concatenation with unit
$\varepsilon$. Resolution $N/p$ is the partial function
\[
N/\varepsilon = N,\qquad
\mathrm{Node}(o, m)/(k \cdot p) = m(k)/p \ \ (k \in \dom m),\qquad
\text{undefined otherwise.}
\]
\end{definition}

\begin{proposition}[Partial monoid action, respected by identity]\label{prop:action}
(i) $N/(p \cdot q) \simeq (N/p)/q$ (Kleene equality: both sides defined together and
equal when defined). (ii) If $N \eqS M$ then for every $p$: $N/p$ is defined iff
$M/p$ is, and $N/p \eqS M/p$ when defined.
\end{proposition}

\begin{proof}
(i) Induction on $p$. (ii) Definedness of $N/p$ is $p \in \Paths(N)$, invariant by
Thm.~\ref{thm:pointing}(a); the value clause follows by induction on $p$ from
Definition~\ref{def:lift}'s pointwise condition.
\end{proof}

Resolution adds no values, constructs nothing on a miss (no autovivification, no
default: whatever $T$ holds, including a null-like payload, is the own value of a node
that exists, never a stand-in for a node that does not), never enters an own value (it walks children, and only children), and
defines no equality. It is a reading of structure
that is already there.

\section{Composition}\label{sec:composition}

A node is exactly its parts, and every existing node can gain a child without anything
else changing. This section states the laws that make both precise. They are the reason
the floor is a product rather than the previous sum (Remark~\ref{rem:previous}). The
operations named here are mathematical descriptions, not required methods of an
implementation.

\begin{proposition}[Complete parts]\label{prop:parts}
Let
\[
\begin{aligned}
\mathrm{decompose}_T &: \NodeT{T} \longrightarrow
  T \times (\finmap{\Bytes}{\NodeT{T}}),\\
\mathrm{compose}_T &: T \times (\finmap{\Bytes}{\NodeT{T}})
  \longrightarrow \NodeT{T},
\end{aligned}
\]
with
\[
\mathrm{decompose}(\mathrm{Node}(o, m)) = (o, m), \qquad
\mathrm{compose}((o, m)) = \mathrm{Node}(o, m).
\]
Then $\mathrm{compose} \circ \mathrm{decompose} = \mathrm{id}$ and
$\mathrm{decompose} \circ \mathrm{compose} = \mathrm{id}$.
\end{proposition}

\begin{proof}
$\mathrm{compose}$ is the structure map of the initial algebra, which is an isomorphism by
Lambek's lemma, and $\mathrm{decompose}$ is its inverse.
\end{proof}

The child map holds whole subtrees, not their own values, and holds every child under its
exact key. A parent's own value is an \emph{input} to reconstruction and is never inferred
from its children: $\mathrm{Node}(o, m)$ and $\mathrm{Node}(o', m)$ with $o \not\approx o'$
have identical children and are different nodes. So a new parent needs a supplied value:
for a generic $T$ nothing can be chosen, while $\NodeT{\mathrm{Option}(T)}$ may supply
$\mathrm{None}$ and $\NodeT{\mathbf{1}}$ the unit (Remark~\ref{rem:instances}).

\begin{definition}[Contexts, split and plug]\label{def:context}
Contexts at a path are defined by recursion on the path. The only context at
$\varepsilon$ is the hole $\square$. A context at $k \cdot p$ is a triple $(o, m, C')$ of an
own value, a finite map with $k \notin \dom m$, and a context $C'$ at $p$. Plugging a node
$s$ into a context gives a node:
\[
\mathrm{plug}(\square, s) = s, \qquad
\mathrm{plug}((o, m, C'), s) = \mathrm{Node}(o,\; m \cup \{k \mapsto \mathrm{plug}(C', s)\}).
\]
Splitting is the partial inverse: $\mathrm{split}_{\varepsilon}(N) = (\square, N)$, and
$\mathrm{split}_{k \cdot p}(\mathrm{Node}(o, n)) = ((o,\, n \setminus k,\, C'),\, s)$ where
$(C', s) = \mathrm{split}_p(n(k))$ and $n \setminus k$ is $n$ with the key $k$ removed; it is
defined when $k \in \dom n$ and the recursive split is.
A context retains every ancestor's own value, the keys leading to the hole, and every
sibling subtree. Writing $\mathrm{Ctx}_T(p)$ for these contexts makes the types explicit:
\[
\begin{aligned}
\mathrm{split}_p &: \NodeT{T} \rightharpoonup
  \mathrm{Ctx}_T(p) \times \NodeT{T},\\
\mathrm{plug}_p &: \mathrm{Ctx}_T(p) \times \NodeT{T} \longrightarrow \NodeT{T}.
\end{aligned}
\]
The path index supplies the key $k$ in each recursive context layer. The domain of
$\mathrm{split}_p$ is exactly the nodes in which $p$ exists.
\end{definition}

\begin{proposition}[Split and plug are inverse]\label{prop:splitplug}
For every $N$ and every $p \in \Paths(N)$, and for every context $C$ at $p$ and every
node $s$,
\[
\mathrm{plug}(\mathrm{split}_p(N)) = N, \qquad
\mathrm{split}_p(\mathrm{plug}(C, s)) = (C, s).
\]
Consequently, writing $N[p \mapsto s]$ for $\mathrm{plug}(C, s)$ with $C$ the context of
$N$ at $p$, which agrees with Definition~\ref{def:replace}:
\[
N[p \mapsto s]/p = s, \qquad N[p \mapsto N/p] = N, \qquad
N[p \mapsto s][p \mapsto t] = N[p \mapsto t].
\]
\end{proposition}

\begin{proof}
Both inverse laws by induction on $p$; at $\varepsilon$ both are the identity. At
$k \cdot p$, plugging the split of $\mathrm{Node}(o, n)$ gives
$\mathrm{Node}(o, (n \setminus k) \cup \{k \mapsto \mathrm{plug}(\mathrm{split}_p(n(k)))\})
= \mathrm{Node}(o, n)$ by the inductive hypothesis. Conversely,
$\mathrm{plug}((o, m, C'), s)$ has $k$ in its domain, removing $k$ gives back $m$ because
$k \notin \dom m$, and the inductive hypothesis gives back $(C', s)$. The replacement laws
follow: resolution of $\mathrm{plug}(C, s)$ at $p$ is $s$ by induction on $p$;
$N[p \mapsto N/p] = \mathrm{plug}(\mathrm{split}_p(N)) = N$; and by the converse law the
context of $N[p \mapsto s]$ at $p$ is again $C$, so replacing there gives
$\mathrm{plug}(C, t) = N[p \mapsto t]$.
\end{proof}

The converse law is the one that is easy to omit. Recovering the context from a result must
return \emph{the same} context, not merely some context that could have produced the
result. The same argument, applied at each path of a finite set $F \subseteq \Paths(N)$ of
which no member is a prefix of another, gives complete cuts. Let $\mathrm{Skel}_T(F)$
retain the ancestor values, keys and all subtrees outside exactly the holes at $F$;
let $\NodeT{T}^{F}$ denote maps with domain exactly $F$ into whole subtrees. Then
\[
\begin{aligned}
\mathrm{cut}_F &: \NodeT{T} \rightharpoonup
  \mathrm{Skel}_T(F) \times \NodeT{T}^{F},\\
\mathrm{rebuild}_F &: \mathrm{Skel}_T(F) \times \NodeT{T}^{F}
  \longrightarrow \NodeT{T}.
\end{aligned}
\]
Cutting is defined exactly when all paths in $F$ exist. These maps are inverse on
their stated domains, including the empty cut and the root-only cut. Different
complete cuts of one tree rebuild the same tree. The holes lie in disjoint
subtrees, so the order in which they are cut does not matter.

\begin{definition}[Containment]\label{def:contain}
$A \sqsubseteq B$ iff $\Paths(A) \subseteq \Paths(B)$ and
$\mathrm{val}_A(p) \approx \mathrm{val}_B(p)$ for every $p \in \Paths(A)$.
\end{definition}

Containment preserves positions, edges and values, not merely a set of values with its
topology forgotten. It is reflexive and transitive, and $A \sqsubseteq B \sqsubseteq A$
gives equal path sets and related values, so $A \eqS B$ by
Theorem~\ref{thm:pointing}: $\sqsubseteq$ is a partial order up to $\eqS$. It needs no order
on $T$ beyond $\approx$. At $T = \mathrm{Option}(A)$ with tag-respecting equality it keeps an
old $\mathrm{None}$ as $\mathrm{None}$; an order in which $\mathrm{None}$ may gain a value is a
different relation, which belongs to the optional interpretation.

\begin{proposition}[Compatible overlay]\label{prop:overlay}
Assume $\approx$ is an equivalence. Trees $A$ and $B$ have a common upper bound
under $\sqsubseteq$ iff
\[
\forall p \in \Paths(A) \cap \Paths(B),\quad
\mathrm{val}_A(p) \approx \mathrm{val}_B(p).
\]
For such a pair their partial overlay $A \sqcup B$ is the least upper bound,
with path set $\Paths(A) \cup \Paths(B)$ and the respective values at those paths,
up to $\approx$. Overlay is commutative, idempotent and associative up to
$\eqS$, with equal definedness for the two associativity expressions.
\end{proposition}

\begin{proof}
A common upper bound relates both values at any shared path to its value there;
symmetry and transitivity give the compatibility condition. Conversely, the union
of the two finite nonempty prefix-closed path sets is again such a set. Assign the
value from $A$ where its path exists, and from $B$ elsewhere. This assignment and
path set determine a finite tree, by grouping nonempty paths at their first key.
Compatibility makes it an upper bound of both trees. Every common upper bound
contains these paths and related values, so it contains their union; this proves
leastness. Commutativity and idempotence follow from the union description.
Either association of three operands is defined exactly when all three pairs
agree on their common paths; when defined, both have the same union of paths and
equivalent values. This proves associativity and equal definedness.
\end{proof}

There is no universal root-only identity for arbitrary $T$: every upper bound must
preserve the root value. $\mathrm{Node}(t,\emptyset)$ is neutral within the component
of trees with root equivalent to $t$. A global least tree exists only when
$T/\!\approx$ has exactly one class. In an optional interpretation the different
order $\mathrm{None}\leq\mathrm{Some}(t)$ may supply an absence-filling overlay;
it is not the relation just proved. Overlay also loses operand provenance:
for $e=\mathrm{Node}(t,\emptyset)$ and $s=\mathrm{Node}(t,\{k\mapsto e\})$,
$e\sqcup s = s\sqcup e = s\sqcup s$ up to identity. The inverse laws for complete
parts do not recover operands of an overlay, edit history, or external effects.

\begin{theorem}[Exact fresh attachment]\label{thm:attach}
For nodes $A, B$, a path $p$ and a key $k$, write $r = p \cdot k$. The operation
$\mathrm{attach}(A, p, k, B)$ has as its domain exactly the arguments with $p \in \Paths(A)$
and $r \notin \Paths(A)$; there, with $A/p = \mathrm{Node}(o, m)$,
\[
\mathrm{attach}(A, p, k, B) \;=\; A[p \mapsto \mathrm{Node}(o,\; m \cup \{k \mapsto B\})].
\]
The right-hand side is defined on the whole domain (availability), and
$C = \mathrm{attach}(A, p, k, B)$ satisfies, exactly,
\[
\Paths(C) = \Paths(A) \cup \{r \cdot q : q \in \Paths(B)\},
\]
\[
\mathrm{val}_C(p') = \mathrm{val}_A(p') \ \ (p' \in \Paths(A)), \qquad
\mathrm{val}_C(r \cdot q) = \mathrm{val}_B(q).
\]
Outside the domain, that is when $p \notin \Paths(A)$ or $r \in \Paths(A)$, it is undefined.
\end{theorem}

\begin{proof}
Definedness: $A/p$ is defined because $p \in \Paths(A)$, and $k \notin \dom m$ because
$r \notin \Paths(A)$, so $m \cup \{k \mapsto B\}$ is a finite partial map and replacement at
$p$ is defined. Undefinedness outside the domain is by definition, not by accident of the
formula: an occupied key is refused even when $B = m(k)$, where the union would still be
a function. Attachment is not an upsert. Paths and
values: replacement keeps the own value and the other children of every node above $p$
(Definition~\ref{def:replace}), and the new node at $p$ keeps $o$ and $m$ and adds exactly
the subtree $B$ under $k$. The new paths $r \cdot q$ cannot coincide with old ones, since
an old path extending $r$ would make $r$ old by prefix closure.
\end{proof}

\begin{corollary}[Consequences of attachment]\label{cor:attach}
In the situation of Theorem~\ref{thm:attach}: $A \sqsubseteq C$; $C/r = B$; the own value
at $p$ is unchanged; restricting $C$ to the paths of $A$, keeping those paths and their
own values, gives back $A$, so every old value is retained as it was, at its old path; and
attachments at $p$ under distinct fresh keys
$k \neq j$ commute:
$\mathrm{attach}(\mathrm{attach}(A, p, k, a), p, j, b) =
\mathrm{attach}(\mathrm{attach}(A, p, j, b), p, k, a)$.
\end{corollary}

\begin{proof}
All but the last are read off the exact paths and values in Theorem~\ref{thm:attach}. The
restriction statement is stronger than $A \sqsubseteq C$, which relates old values only up
to $\approx$; attachment keeps them exactly. For commutation, both orders put
$\mathrm{Node}(o, m \cup \{k \mapsto a, j \mapsto b\})$ at $p$ and change nothing else.
\end{proof}

Availability is the half that separates this model from the previous one, and a law
stated only for successful attachments would not capture it: an operation that refused
every childless parent, or that returned $A$ unchanged, would satisfy $A \sqsubseteq C$
vacuously or trivially. Inside the target of $E$ (Remark~\ref{rem:previous}), a child
can be attached below $\mathrm{Node}(\mathrm{Some}(x), \emptyset)$, the image of
$\Leaf(x)$, and the result lies outside the image of $E$. In the previous model the same node could gain a child only by
being replaced, which moved its value to another path.

\begin{proposition}[No path- and value-preserving map to the previous model]\label{prop:nonembed}
Let $T$ be inhabited. Read the previous model's observations through $E$: a previous tree
$L$ has the paths of $E(L)$, and the optional value $\mathrm{val}_{E(L)}(p)$ at each of them.
Then there is no function $\Phi : \NodeT{T} \to L(T)$ with
\[
\Paths(\Phi(N)) = \Paths(N) \quad\text{and}\quad
\mathrm{val}_{E(\Phi(N))} = \mathrm{Some} \circ \mathrm{val}_N
\]
for every $N$.
\end{proposition}

\begin{proof}
Take $x \in T$ and $N = \mathrm{Node}(x, \{k \mapsto \mathrm{Node}(x, \emptyset)\})$. If
$\Phi$ existed, $E(\Phi(N))$ would have the value $\mathrm{Some}(x)$ at $\varepsilon$ and the
path set $\{\varepsilon, k\}$. A previous tree with
a value at $\varepsilon$ is $\Leaf(x')$ for some $x'$, and $E(\Leaf(x')) =
\mathrm{Node}(\mathrm{Some}(x'), \emptyset)$ has the path set $\{\varepsilon\}$. So no
previous tree has both observations, and $\Phi(N)$ cannot exist.
\end{proof}

The statement is about the required interface, not about bijections between underlying
sets: a lossless profile of this model inside the previous one exists, at the cost of
changed native paths. The design records give this argument informally; the proposition
makes it formal, and slightly stronger, since it rules out every function
with those properties and not only isomorphisms.

\section{Slots: admission and its tiers}\label{sec:slots}

Every requirement deixis places on a slot is a requirement on the pair $(T, \approx)$,
never on the carrier alone. The capabilities form a ladder of tiers --- and, answering
a gap between the specification's motivation (``identity decided by looking'') and its
formal minimum (an equivalence), the effective tiers are now explicit:

\begin{center}
\begin{tabular}{@{}lll@{}}
\toprule
supplied & available & \\
\midrule
carrier $T$ & shape, keys, paths, navigation & \S\ref{sec:tree}--\ref{sec:path}\\
$+\ \approx$ an equivalence & \emph{semantic} identity $\eqS$ & \S\ref{sec:identity}\\
$+$ a decision procedure for $\approx$ & \emph{effective} identity & Prop.~\ref{prop:decidable}\\
$+\ \enc : T \to \Bytes$ lawful, computable & canonical bytes, hashing, addressing & \S\ref{sec:codec}\\
$+$ an exact partial inverse for $\enc$ & decodable canonical representation & \S\ref{sec:codec}\\
\bottomrule
\end{tabular}
\end{center}

\noindent where \emph{lawful} means the complete-invariant law
\begin{equation}\label{eq:codeclaw}
x \approx y \iff \enc(x) = \enc(y),
\end{equation}
i.e.\ $\enc$ factors through the quotient $T/\!\approx$ and is injective on it. One
inclusion between the tiers is a theorem: a computable lawful $\enc$ automatically
supplies a decision procedure for $\approx$ (compute and compare), so the
canonical-bytes tier contains the effective tier. We claim \emph{no} strictness in
the other direction. Prop.~\ref{prop:gap} separates a different pair of notions ---
set-theoretic existence from computable existence, over a slot whose $\approx$ is
undecidable --- and says nothing about whether decidable equality falls short of a
computable invariant. Indeed, in natural effective presentations the gap can close
outright: on $\mathbb{N}$ with any decidable equivalence $E$, the least-representative
map $c(n) = \min\{\, m \le n : m \mathrel{E} n \,\}$ is a computable complete
invariant. Whether some fixed effective model makes the separation strict is a
question about that model's definition of ``computable on $T$,'' and this paper does
not pose one. A slot admitted only at the semantic tier remains mathematically
legitimate and operationally mute: its node identity is well-defined and
undecidable.

\begin{theorem}[Existence boundary, non-effective]\label{thm:countable}
A function $\enc$ satisfying \eqref{eq:codeclaw} exists (as a set-theoretic function,
with no computability claim) iff $T/\!\approx$ is countable.
\end{theorem}

\begin{proof}
$\Bytes$ is countably infinite. A lawful $\enc$ induces an injection
$T/\!\approx\ \hookrightarrow \Bytes$, forcing countability; conversely any injection
of the quotient into $\Bytes$, composed with the projection, is lawful.
\end{proof}

\begin{proposition}[Existence and computability come apart]\label{prop:gap}
Let $T \subseteq \mathbb{N}$ be the set of indices of total computable functions
$\mathbb{N} \to \mathbb{N}$ and $i \approx j$ iff $\varphi_i = \varphi_j$
extensionally. Then $T/\!\approx$ is countable, so a lawful $\enc$ exists
(Thm.~\ref{thm:countable}); but there is no computable procedure that halts on every
$i \in T$ and is lawful there --- neither a total computable function on all indices
whose restriction to $T$ is lawful, nor a promise procedure guaranteed to halt on the
promise $T$. (Since $T$ is not decidable, ``computable on $T$'' must be one of these
two readings; the proof covers both.)
\end{proposition}

\begin{proof}
Countability is immediate. Suppose $\enc$ is computable in either reading and lawful
on $T$. Fix an index $z \in T$ of the constant-zero function. Given any program $e$
and input $x$, the s-m-n theorem \cite{soare} effectively yields an index $g(e,x)$ of
\[
f_{e,x}(n) =
\begin{cases}
1, & \text{if } e(x) \text{ halts within } n \text{ steps},\\
0, & \text{otherwise,}
\end{cases}
\]
which is total for \emph{every} $e, x$ --- so $g(e,x) \in T$ always, and the promise
is met. Now $f_{e,x} = \varphi_z$ (extensionally) iff $e(x)$ never halts. Computing
and comparing $\enc(g(e,x))$ and $\enc(z)$ --- both calls halt, since both indices
satisfy the promise --- decides, by lawfulness on $T$, whether $g(e,x) \approx z$,
hence whether $e(x)$ halts. That decides the halting problem; contradiction.
\end{proof}

\begin{proposition}[Decidability transports]\label{prop:decidable}
$\eqS$ is decidable iff $\approx$ is. (From $\approx$ to $\eqS$ by structural recursion:
key comparison and domain comparison are decidable outright, and own values are compared
by $\approx$. Conversely, $x \approx y$ iff
$\mathrm{Node}(x, \emptyset) \eqS \mathrm{Node}(y, \emptyset)$.)
\end{proposition}

Two consequences frame the design space. Since every carrier admits at least the
discrete and the total relation, \emph{no carrier is inadmissible}; only relations are
(a non-equivalence lifts to a non-equivalence,
Thm.~\ref{thm:equiv}). And since whatever $\approx$ merges is merged for every layer
above --- Prop.~\ref{prop:action}(ii) shows no \emph{extensional} observation through
the structure separates merged values, and Remark~\ref{rem:extensional}'s admission
classes ensure every admitted map is at least extensional --- the choice of $\approx$
is unrecoverable in the disciplined ecosystem, which is why the specification
counsels choosing it as fine as any consumer will need.

\section{Canonical bytes: the codec combinator}\label{sec:codec}

deixis owns only half an encoding: the structural half (tags, counts, key bytes,
entry order --- discharged by $\Bytes$ and $\lex$); the slot half is the supplied
$\enc$. What the floor can define is a combinator.

\begin{definition}[Framing assumptions]\label{def:framing}
Fix a length code $\langle \cdot \rangle : \mathbb{N} \to \Bytes$ that is
\emph{injective}, \emph{prefix-free}, and equipped with a unique exact partial
decoder --- for every byte stream, at most one prefix is in the code's image, and the
decoder recovers its unique preimage and rejects everything else, including
non-canonical spellings of a length (e.g.\ shortest-form unsigned LEB128 with
non-shortest forms rejected) --- and a fixed marker string $\tau$, which may be empty.
\end{definition}

\begin{definition}[The combinator]\label{def:encd}
Given lawful $\enc$, define $\encD : \NodeT{T} \to \Bytes$ by structural recursion:
\[
\begin{array}{lcl}
\encD(\mathrm{Node}(o,m)) &=&
  \tau \cdot \langle |\enc(o)| \rangle \cdot \enc(o) \cdot \langle |m| \rangle\\[2pt]
&& {}\cdot
  \displaystyle\prod_{k \in \dom m}^{\lex\text{-ascending}}
  \langle |k| \rangle \cdot k \cdot \encD(m(k))
\end{array}
\]
where the product concatenates entries in strictly $\lex$-ascending key order (keys
are unique, so the order is total and the sequence strict). This is an abstract
encoder family over the mandatory-value model. ADR~0011 of the specification selects
$\tau$ empty for \texttt{deixis-codec-v2}; that codec's contract, not this paper, fixes
the byte layout, identifier, linked form and refusal protocol. Nothing below depends on
those choices, or on the length of $\tau$.
\end{definition}

\begin{lemma}[Self-delimitation]\label{lem:selfdelim}
$\encD$ is a prefix code: for all $N, M$, if $\encD(N)$ is a prefix of $\encD(M)$
then $\encD(N) = \encD(M)$; moreover a byte stream determines its parse as
(tag, framed fields) uniquely.
\end{lemma}

\begin{proof}
By induction on nodes, using prefix-freeness of $\langle\cdot\rangle$. After the
fixed marker, the own-value length determines that payload's extent. The child
count fixes how many entries follow; each key-length field fixes its key's extent,
and the inductive hypothesis fixes its child's extent. This determines a unique
end for the whole node, including one with no children, so no complete encoding
can be a proper prefix of another.
\end{proof}

\begin{theorem}[Lifted codec law]\label{thm:codec}
Let $\enc$ be lawful for the slot setoid $(T,\approx)$.
Its lifted encoder is lawful for nodes:
\[
N \eqS M \iff \encD(N) = \encD(M).
\]
\end{theorem}

\begin{proof}
($\Rightarrow$) Induction. Related own values encode equally by
\eqref{eq:codeclaw}. Equal domains yield the identical $\lex$-sorted key sequence;
pointwise $\eqS$ children encode equally by the inductive hypothesis; the
concatenations coincide.
($\Leftarrow$) Equal byte strings parse identically (Lemma~\ref{lem:selfdelim}):
equal own-value payloads give related values by \eqref{eq:codeclaw};
equal counts, then entry-by-entry equal keys and equal child
encodings, giving equal domains and, by the inductive hypothesis, pointwise
$\eqS$ children --- exactly Definition~\ref{def:lift}.
\end{proof}

\begin{corollary}[Content addressing, computational]
Let $H$ be a digest. Unconditionally, $N \eqS M \Rightarrow H(\encD(N)) =
H(\encD(M))$: equal nodes always share an address. The converse is not a mathematical
fact --- no fixed-size digest is injective on an infinite domain --- but a
computational one: under a collision-resistance assumption on $H$, exhibiting
$N \not\eqS M$ with $H(\encD(N)) = H(\encD(M))$ is infeasible. Address equality is
therefore evidence, not proof, of node identity: retrieved content must be rehashed
against its address, and where identity is load-bearing, its canonical bytes
revalidated. With that discipline, hash-keyed leaves recover shared substructure and
DAGs without a cycle ever entering a value.
\end{corollary}

\begin{remark}[Encoder versus codec]\label{rem:codecword}
What Theorem~\ref{thm:codec} delivers is a computable \emph{complete invariant} ---
an encoder. A \emph{codec} additionally ships the exact partial decoder: a normative
parser that reconstructs the node \emph{up to identity} and rejects, with stable
reasons, every byte string outside the canonical image --- malformed or non-canonical
lengths, malformed slot-codec identifiers, unsorted or duplicate keys, counts
inconsistent with the stream, payloads outside the slot codec's image, trailing bytes,
and resource-limit breaches (depth, size) as implementation policy distinct from the
value set. The decoder laws are setoid-level, because a
lawful $\enc$ is injective only on the quotient: $D(\encD(N)) \eqS N$ (never raw
equality of syntax trees --- a slot relation coarser than representation makes raw
recovery impossible by design), and $\encD(D(b)) = b$ byte-exactly on the decoder's
domain. Countability (Thm.~\ref{thm:countable}) bounds the existence of
\emph{invariants}; an effective codec needs the decodable tier of \S\ref{sec:slots}.
The specification deliberately freezes no byte layout yet. Its candidate,
\texttt{deixis-codec-v2} (\texttt{CODEC.md}, ADR~0011), instantiates this model with
one node production --- the node's own payload, length-framed, then its count and
entries --- and $\tau$ empty. Its framing and identity laws are proved separately, in
prose, against that concrete grammar (\texttt{CODEC-proofs.md}), and are pinned by
independently authored vectors. The earlier leaf/structure grammar survives only as the
withdrawn \texttt{deixis-codec-v1}. Until the freeze (a clean-room implementation and
a signed manifest) the grammar remains a candidate; this theorem does not discharge
that delivery gate. A \emph{deployed} content address must additionally
bind the codec version, the slot codec identifier, and the digest algorithm (domain
separation or a typed envelope) --- an interoperability obligation of the frozen
artifact, outside this paper's scope.
\end{remark}

\section{The positional spelling}\label{sec:pos}

Sequence semantics enter deixis by writing positions into keys. The profile
\texttt{deixis-pos-v1} fixes the spelling once, generically.

\begin{definition}[$\magn$ and $\kappa$]\label{def:kappa}
For $i \in \mathbb{N}$ let $\magn(i)$ be the shortest unsigned big-endian encoding of
$i$, except $\magn(0) = \mathtt{00}$ (one octet; the only magnitude with a leading
zero octet). With $r(i) = |\magn(i)| - 1$:
\[
\kappa(i) \;=\; \mathtt{ff}^{\,r(i)} \cdot \mathtt{00} \cdot \magn(i),
\qquad |\kappa(i)| = 2\,|\magn(i)|.
\]
\end{definition}

The spelling has an honest lineage: $\kappa$ is a byte-alphabet Elias-gamma code
\cite{elias} with the unary length prefix written in the \emph{maximal} symbol.
Gamma over bits ($0^{k}1 \cdot \mathrm{bin}(i)$) is prefix-free but
order-\emph{reversing} under lexicographic comparison, since leading zeros sort low;
writing the run in $\mathtt{ff}$ is exactly the inversion that makes length dominate
in the right direction. The construction belongs to the ordered-key family; the
contribution here is its ratification as a named, versioned profile with the three
laws proved and vector-pinned.

\begin{theorem}[The three $\kappa$ laws]\label{thm:kappa}
\emph{(i)} $\kappa$ is injective.
\emph{(ii)} $\kappa$ is prefix-free: no $\kappa(i)$ is a proper prefix of a
$\kappa(j)$.
\emph{(iii)} $\kappa$ is order-preserving: $i < j \iff \kappa(i) \lex \kappa(j)$.
\end{theorem}

\begin{proof}
Write $r_i = r(i)$. Observe first that the $\mathtt{ff}$-run length of $\kappa(i)$ is
exactly $r_i$: the run is terminated by the $\mathtt{00}$ octet, and $\magn(i)$'s
leading octet is $\mathtt{ff}$ only when it would extend the run --- but then the
parse below still recovers $r_i$ from the terminator position, which is octet
index $r_i$ by construction.\footnote{Concretely: octets $0,\dots,r_i-1$ are
$\mathtt{ff}$ and octet $r_i$ is $\mathtt{00}$; whether later octets are
$\mathtt{ff}$ is irrelevant to the claims, which only compare octets up to index
$\max(r_i, r_j)$ and, on equal $r$, the magnitudes.}

(ii) Suppose $r_i < r_j$. At octet index $r_i$, $\kappa(i)$ carries $\mathtt{00}$
while $\kappa(j)$ carries $\mathtt{ff}$ (still inside $j$'s run). They differ at an
index interior to both, so neither is a prefix of the other. If $r_i = r_j$, the two
have equal length $2(r_i{+}1)$, and a prefix relation forces equality.

(i) If $\kappa(i) = \kappa(j)$ then $r_i = r_j$ (terminator position) and
$\magn(i) = \magn(j)$; shortest-form magnitudes of equal length are equal iff the
numbers are, and the $\magn(0)$ exception is the unique single-octet magnitude with
leading $\mathtt{00}$.

(iii) If $r_i < r_j$: numerically, shortest form gives
$i < 256^{\,r_i+1} \le 256^{\,r_j} \le j$ (the middle inequality since
$r_i + 1 \le r_j$; the last since $|\magn(j)| = r_j{+}1$ forces $j \ge 256^{r_j}$,
noting $r_j \ge 1$ here); lexicographically, the first differing octet is index
$r_i$, where $\mathtt{00} < \mathtt{ff}$. So larger run implies both larger number
and lexicographically larger key. If $r_i = r_j$: the prefixes agree through the
terminator, and equal-length big-endian magnitude comparison is numeric comparison.
Together with (i), trichotomy on both sides yields the equivalence.
\end{proof}

\begin{corollary}[Sequences at zero interpretive cost]\label{cor:pos}
For a node whose child domain is
$\{\kappa(0), \dots, \kappa(n{-}1)\}$, canonical
($\lex$-ascending) entry order equals index order. Consequently: a consumer walking a
sequence generates $\kappa(0), \kappa(1), \dots$ and compares octets --- no key is
ever parsed, and no decoder need exist. Membership of an arbitrary byte string in
$\mathrm{im}\,\kappa$ is decidable by the run/terminator/shortest-form checks alone.
\end{corollary}

The specification's recognition discipline --- \emph{exact or refusal} --- is a
corollary of the profile laws, not of extensionality alone: a reader that
``repairs'' a non-canonical spelling (e.g.\ reads $\mathtt{ff000001}$ as position
$1$) is a recognizer that identifies two nodes $\eqS$ distinguishes --- it coarsens
identity \emph{through the reading}, breaking exactness and identity transport
(Def.~\ref{def:profile}(ii),(iii)). Extensionality would not forbid it (a repairing
reader can be perfectly extensional); the identity-bearing obligations of
Remark~\ref{rem:extensional}(b),(c) do --- and once such a reading is deployed, no
extensional observer downstream could detect or undo the coarsening, which is why the
refusal is placed at the recognizer.

\begin{remark}[Where prefix-freeness earns its keep]\label{rem:prefixfree}
The floor's path type never concatenates keys --- a path is a \emph{sequence} of
already-delimited keys, and resolution consumes them one at a time --- so the
sequence embedding below uses only injectivity and order preservation. Prefix-freeness
is provided for the \emph{flattened} contexts profiles may define: spelling a whole
key path as a single byte string (single-string path names, radix layouts, streamed
prefixes), where unambiguous parsing of concatenated keys is exactly the prefix-free
property. It costs $\kappa$ nothing and licenses those readings in advance.
\end{remark}

\section{Structure profiles: container semantics as embeddings}\label{sec:profiles}

The branching functor of Def.~\ref{def:functor} is one point in the space of
containers; abstractly one could parameterize it,
$\mathrm{Node}(C, T) = \mu X.\ T \times C(X)$, obligating each $C$ to a lifted
congruence and a canonicalization (a \emph{setoid functor}). deixis refuses the
parameter: distinct instances would be distinct floors with incomparable identities,
and the floor's value is its uniqueness. Instead:

\begin{definition}[Structure profile]\label{def:profile}
Let $(A, =_A)$ be the carrier of a reading's semantics (a setoid). A
\emph{structure profile} for $A$ over the slot $(T, \approx)$ is a named, versioned
pair
\[
P : A \to \NodeT{T}, \qquad R : \NodeT{T} \rightharpoonup A
\]
of extensional maps satisfying:
\begin{enumerate}[label=(\roman*), leftmargin=2.2em]
\item \emph{retraction}: $R(P(a)) =_A a$ for all $a$;
\item \emph{exactness}: if $R(n)$ is defined then $P(R(n)) \eqS n$ --- recognition
  inverts onto the same node, so $\mathrm{dom}\,R$ is, up to $\eqS$, exactly
  $\mathrm{im}\,P$;
\item \emph{identity transport}: $a =_A b \iff P(a) \eqS P(b)$ --- the reading's
  identity reduces to node identity, never extends it;
\item \emph{recognition}: $R$ is extensional ($n \eqS n'$ implies equidefinedness and
  $=_A$-equal results) and refuses rather than repairs --- no normalization of a node
  outside $\mathrm{im}\,P$ into it is permitted. Effectiveness of recognition is the
  \emph{profile's own} obligation, stated relative to the slot's tiers and discharged
  per profile (the sequence profile by its dense-domain test, the set profile by its
  coherence assumptions); it does not follow automatically from the slot being
  effective;
\end{enumerate}
together with the spelling laws its container's semantics require of the keys $P$
writes: \emph{injectivity} always; \emph{order preservation} iff the reading carries
order (so canonical entry order restates it); \emph{prefix-freeness} where flattened
spellings are intended (Remark~\ref{rem:prefixfree}). This is the ontos bridge's
$(P, R)$ structure (\S\ref{sec:ontos}), promoted to the general interface.
\end{definition}

A profile declares its carrier. A container node needs an own value like every node,
and the floor supplies none (Remark~\ref{rem:instances}). The two profiles below take
the slot to be $\mathrm{Option}(U)$: $\mathrm{None}$ marks the container node, and a node
carrying $\mathrm{Some}$ is refused as a container. This is the representation the set
profile's specification record states and its migrated implementations use.
Another carrier would need its own image and recognizer.

\begin{remark}[Law dependencies; the categorical name]\label{rem:profilelaws}
The four laws are deliberately redundant, and the dependency structure matters for
conformance. Law (iii) is derivable: the forward direction is extensionality of $P$;
backward, $P(a) \eqS P(b)$ gives $a =_A R(P(a)) =_A R(P(b)) =_A b$ by
(iv)-extensionality of $R$ together with (i). The no-repair clause of (iv) restates
(ii). A minimal conformance suite therefore generates (i), (ii), and (iv); testing
(iii) directly remains useful because vectors cannot observe extensionality itself.
Categorically, a profile is a complemented subobject with chosen complement (the
refusals) --- equivalently a lawful prism $\NodeT{T} \cong A \sqcup
\mathrm{Refused}$ on quotients --- and exactness says the idempotent
$e = P \circ R$ is a \emph{restriction} idempotent in the sense of restriction
categories \cite{cockettlack}: not merely $e \circ e = e$ but $e(n) \eqS n$ on its
domain. A repair pass is precisely a split idempotent that is \emph{not} a
restriction idempotent --- a retraction that moves points onto the image --- which is
the one-line categorical form of exact-or-refusal.
\end{remark}

\begin{proposition}[Sequence transport]\label{prop:pos-transport}
Let $T = \mathrm{Option}(U)$ and take $A = \NodeT{T}^{*}$ with pointwise-and-length
equality. Define
$P_{\mathrm{seq}}(N_0, \dots, N_{n-1}) = \mathrm{Node}(\mathrm{None}, \{\kappa(i) \mapsto
N_i\})$ and let $R_{\mathrm{seq}}(\mathrm{Node}(o, m))$, for $n = |\dom m|$, be defined iff
$o = \mathrm{None}$ and $\dom m = \{\kappa(0), \dots, \kappa(n-1)\}$, returning the
children in index order. A node carrying $\mathrm{Some}$ is refused whatever its children:
accepting it would read two different nodes, differing only in their own values, as the
same sequence. Then $(P_{\mathrm{seq}}, R_{\mathrm{seq}})$ satisfies
Definition~\ref{def:profile} (i)--(iv), with identity transport
\[
P_{\mathrm{seq}}(\vec{N}) \eqS P_{\mathrm{seq}}(\vec{M})
\iff |\vec N| = |\vec M| \wedge \forall i.\ N_i \eqS M_i .
\]
\end{proposition}

\begin{proof}
$\kappa$'s injectivity (Thm.~\ref{thm:kappa}(i)) makes $P_{\mathrm{seq}}$ a function
with $n$ distinct keys; the identity equivalence is Definition~\ref{def:lift} for
nodes carrying $\mathrm{None}$ plus index-wise key matching. (i) and (ii) are immediate from the
own-value and dense-domain conditions; (iv): the own-value condition is a tag test and
the domain condition is decided by generate-and-compare (Cor.~\ref{cor:pos}), and
extensionality holds because $\eqS$-equal nodes have related own values (so the same
tag, under the slot's tag-respecting equality), equal domains and pointwise-$\eqS$ children, and no non-dense domain is ever read back into
a sequence. This is exactly the ontos recognizer of \S\ref{sec:ontos}, with children
left as nodes.
\end{proof}

\begin{proposition}[Set transport, repaired]\label{prop:set-transport}
Let $(S, \approx_S)$ be the member sort, $\enc$ lawful for it
\eqref{eq:codeclaw}, and let $\hat\mu : S/\!\approx_S \to \NodeT{T}$ be a
\emph{member-form assignment} --- a function on classes (a chosen member form per
class; ``assignment'' rather than ``section,'' since it is not a right inverse of any
given projection), over the slot $T = \mathrm{Option}(U)$. No injectivity of $\hat\mu$ is
assumed: keys, not children,
separate classes, since $\enc$ is already injective on the quotient. Take
$A = \mathcal{P}_{\mathrm{fin}}(S/\!\approx_S)$ with extensional equality, and (with
$\enc$ constant on classes, written $\enc(c)$):
\[
P_{\mathrm{set}}(A) \;=\; \mathrm{Node}(\mathrm{None},\; \{\, \enc(c) \mapsto \hat\mu(c) \;:\; c \in A \,\}).
\]
This is literally a finite map --- one entry per class, keys distinct by injectivity
of $\enc$ on the quotient. Suppose further the slot supplies \emph{decidable
coherence}: for a key--child pair one can decide whether some (necessarily unique)
class $c$ has $\enc(c) = k$ and $\hat\mu(c) \eqS \text{child}$ --- e.g.\ via an exact
partial decoder for $\enc$ (the decodable tier) plus decidable $\approx_S$ and slot
equality. Let $R_{\mathrm{set}}$ be defined on exactly the nodes carrying $\mathrm{None}$
all of whose entries cohere, returning the set of their classes. A node carrying
$\mathrm{Some}$ is refused, since accepting it would make set identity coarser than node
identity. (Deciding coherence needs
effective access to the assignment itself --- computing $\hat\mu(c)$ from a decoded
member --- not only a decoder for the key invariant. At the decodable tier no choice
is involved at all: take $\hat\mu(c) := \text{member-form}(D(\enc(c)))$ --- the
exact partial decoder \emph{is} the canonical assignment.) Then
$(P_{\mathrm{set}}, R_{\mathrm{set}})$ satisfies Definition~\ref{def:profile}
(i)--(iv); membership transports ($c \in A \iff \enc(c) \in \dom$); and insertion is
idempotent by construction (a repeated member spells a duplicate key, which the floor
refuses to express).
\end{proposition}

\begin{proof}
$P_{\mathrm{set}}$ is well-defined \emph{as a function into raw nodes} precisely
because $\hat\mu$ is a function on classes --- this is the repair: with a member form
chosen per class, no ambiguity of raw children arises. (Injectivity of $\hat\mu$
plays no role in any clause below.) (iii): equal node identity gives equal
domains, hence equal class sets by injectivity of $\enc$ on the quotient; conversely
equal sets give identical maps. (i): each entry of $P_{\mathrm{set}}(A)$ coheres by
construction and $R_{\mathrm{set}}$ reads its classes back. (ii): if $n$ carries $\mathrm{None}$ and every
entry coheres then $n$ and $P_{\mathrm{set}}(R_{\mathrm{set}}(n))$ both carry $\mathrm{None}$,
have equal domains and $\eqS$-equal children (coherence pins each child to $\hat\mu(c)$ up to $\eqS$).
(iv): coherence is decidable by hypothesis and the own-value condition is a tag test;
extensionality follows since coherence respects $\eqS$ of children and $\eqS$-equal
nodes carry the same tag; nothing outside the coherent image is normalized in.
\end{proof}

\begin{remark}[The setoid-valued variant]
The shipping specification of the set profile deliberately does \emph{not} choose
representatives: any $\approx_S$-equal member may sit under the key
(``representatives are free''), recognition checks coherence entry by entry, and all
laws hold with $P_{\mathrm{set}}$ valued in $\NodeT{T}/\!\eqS$ rather than in raw
nodes. The two presentations agree up to $\eqS$; the assignment is what an implementation
that must \emph{produce} a node uses, and the quotient-valued reading is what
recognition-first conformance pins. Stating the raw-node version without a chosen
assignment --- as an earlier draft of this paper did --- is a genuine error: distinct
representatives yield $\eqS$-equal but unequal nodes, and the comprehension fails to
define a map at all.
\end{remark}

\begin{remark}[Representation scale]\label{rem:setscale}
Proposition~\ref{prop:set-transport} is the \emph{self-keyed} presentation, and its
laws are exact at every member sort --- but as a physical representation it
duplicates each member's complete invariant in the key beside the child. For slot
members with bounded canonical bytes this is a compact spelling; for \emph{node}
members under $\enc = \encD$ it compounds: the key of a set containing $S$ embeds
all of $\encD(S)$, so singleton nestings $S_{i+1} = \{S_i\}$ double in size,
and no linked representation rescues this --- the parent \emph{key} still spells the
member's complete flat encoding. The duplication is intrinsic to exact direct
self-keying (an injective key for an arbitrary member must carry enough bytes to
distinguish it), not to the transport laws: the same finite sets embed lawfully as a
\emph{sorted positional} spelling --- dense $\kappa$ keys with children strictly
increasing under a profile-owned total order induced by a complete invariant of
member identity, one occurrence per class enforced by strictness --- trading
one-lookup membership for search. Which representation a profile freezes is an
engineering choice the laws do not make; the shipping specification ratifies
self-keying for slot members and, for node members, supersedes direct self-keying
in favor of the sorted positional shape, gated on the frozen codec.
\end{remark}

\begin{remark}[Scope of the principle]\label{rem:scope}
Multiplicity (count-valued maps or a positional embedding) and records (authored name
keys) are further profiles of the same shape. But the principle as \emph{proved}
extends exactly this far: positional and collection readings --- sequences, sets,
their composites --- enter as profiles; no universality theorem for arbitrary
container semantics is claimed. Two obstacles delimit any general statement. A
container has \emph{shape} as well as positions: two distinct nullary shapes would
both spell the empty map unless a shape marker (a reserved key namespace) is part of
the spelling. And the key space is countable, so only countably presented finitary
containers can embed uniformly. We conjecture, and leave as future work, the scoped
theorem: every countably presented finitary container with decidable shape codes
embeds via a shape marker plus per-shape position spellings in disjoint key
namespaces; quotient-symmetric collections (sets, bags) are not instances of the raw
polynomial frame at all --- they are \emph{quotient containers} \cite{aagm} --- and
enter, as above, through a complete invariant of member identity where a lawful
$\enc$ exists --- as canonical keys (Prop.~\ref{prop:set-transport}) or as a
canonical order (Rem.~\ref{rem:setscale}): an
orbit-splitting by a complete invariant, in that vocabulary. Within that scope the moral stands: openness lives in the two
dimensions identity does not interpret --- the key space and the slot --- readings
compose in one node under one identity, where parameters would stack functors and
fragment the floor.
\end{remark}

\begin{remark}[Signaling and composition --- proposed mechanism]\label{rem:signal}
How does a consumer learn \emph{which} profile to apply? The floor's commitment is
precise: deixis does not force a profile tag into every payload node --- a mandatory
in-node tag would leak semantics into identity (two shape-identical payloads with
different tags would be different nodes) and would invite exactly the
tag-versus-shape repairs the discipline forbids. A node may lie in several profile
images at once, and deliberately so --- overlap is classification, not ambiguity of
identity: a dense-$\kappa$ node carrying $\mathrm{None}$ \emph{is} a sequence and a degenerate
record, and
generic tooling exploits that. Above the floor, two distinct mechanisms are
\emph{proposed} (design in review; labeled as such, not counted among this paper's
results). \emph{Out-of-band claim}: a profile descriptor --- named, versioned,
hashable --- travels beside the payload; it does not alter the payload's node
identity, has its own identity, and is checked by running the named recognizer.
\emph{Self-describing wrapper}: a wrapper profile whose reserved key carries the
payload's descriptor by hash, with the payload under another key --- the descriptor
then \emph{does} enter the wrapper's structural identity, deliberately, while the
enclosed payload retains its own. Descriptor fields, canonical encoding, payload
binding, and behavior under an unavailable profile implementation are obligations of
that mechanism's own record. Composition at one structural level is governed by the
spelling laws: each profile writes into a declared key region (dense $\kappa$-keys;
invariant keys; authored-name keys under a shape marker per Remark~\ref{rem:scope}),
composed profiles must declare disjoint regions, and a composition whose declared
regions overlap is refused as ill-formed rather than resolved by precedence. Version
and digest binding for \emph{deployed} claims is the frozen codec artifact's
obligation (Remark~\ref{rem:codecword}), outside this paper's scope.
\end{remark}

\section{The ontos instance}\label{sec:ontos}

ontos is a value model in production use: $\VO = \mu X.\ \Bytes \sqcup X^{*}$, with
constructors $\Atom$ and $\Tup$, structural identity $\eqO$, and a frozen canonical
codec. It predates deixis; the claim proved here is that it is, exactly, a
dense-position profile of $\NodeT{\mathrm{Option}(\Bytes)}$. It is the earlier
projection followed by $E$ from Remark~\ref{rem:previous}.

Instantiate $T = \mathrm{Option}(\Bytes)$ with tag-respecting octet equality
(so $\eqS$ is absolute).
Define the total projection $P$ and partial recognizer $R$:
\[
\begin{array}{lcl}
P(\Atom\,b) = \mathrm{Node}(\mathrm{Some}(b),\emptyset), &&\\[2pt]
P(\Tup(v_0 \dots v_{n-1})) = \mathrm{Node}(\mathrm{None},\{\kappa(i) \mapsto P(v_i)\}), &&\\[3pt]
R(\mathrm{Node}(\mathrm{Some}(b),\emptyset)) = \Atom\,b, &&\\[2pt]
R(\mathrm{Node}(\mathrm{None},m)) = \Tup(R(m(\kappa(0))), \dots, R(m(\kappa(n{-}1)))), && n = |\dom m|.
\end{array}
\]
The tuple clause is defined only when $\dom m = \{\kappa(0), \dots, \kappa(n{-}1)\}$
and every child recognizes. A node carrying $\mathrm{Some}(b)$ and children refuses.
The complete types are $P : \VO\to\NodeT{\mathrm{Option}(\Bytes)}$ and
$R : \NodeT{\mathrm{Option}(\Bytes)}\rightharpoonup\VO$.

\begin{theorem}[Exactness of the bridge]\label{thm:ontos}
\emph{(i)} $R(P(v)) = v$ for all $v \in \VO$.
\emph{(ii)} $R(d)$ is defined iff $d \in \mathrm{im}\,P$; and $R(d) = v$ implies
$P(v) \eqS d$.
\emph{(iii)} $P$ is an embedding of identities:
$v \eqO w \iff P(v) \eqS P(w)$.
\emph{(iv)} Pointing commutes: the index path $(i_0, \dots, i_k)$ resolves in $v$ iff
the key path $(\kappa(i_0), \dots, \kappa(i_k))$ resolves in $P(v)$, to the
projection of the same subvalue.
\end{theorem}

\begin{proof}
(i) Induction on $v$; the tuple case uses $\kappa$'s injectivity for the domain and
Corollary~\ref{cor:pos} to read the children back in index order.
(ii) ($\Leftarrow$) is (i). ($\Rightarrow$) Induction on $d$: the $\mathrm{None}$ case's
definedness condition forces the domain to be exactly $\{\kappa(0..n{-}1)\}$ and the
children into $\mathrm{im}\,P$ inductively, so $d = P(R(d))$ up to $\eqS$; at
$T = \mathrm{Option}(\Bytes)$ with tag-respecting octet equality, $\eqS$-equal nodes
are equal, giving both claims. The $\mathrm{Some}$ case requires no children and
is the atom clause exactly.
(iii) By (i) and functionality of $R$ on its domain: $P(v) \eqS P(w)$ implies (at
this slot, equality) $v = R(P(v)) = R(P(w)) = w$; the converse is congruence of the
defining recursion.
(iv) Induction on the path, unfolding $P$ and Definition~\ref{def:path}; the step
case is exactly the tuple clause of $P$.
\end{proof}

\begin{corollary}[Frozen bytes survive the retraction]
With $\mathrm{enc_o}$ ontos' frozen codec,
\[
\mathrm{enc_o}(R(P(v))) = \mathrm{enc_o}(v).
\]
This is immediate from
Thm.~\ref{thm:ontos}(i), and worth stating only for what it does \emph{and does not}
say. It says the round trip through deixis is invisible to the incumbent encoding:
project, recognize, and every downstream hash and signature over
$\mathrm{enc_o}$-bytes is untouched. It does \emph{not} say
$\encD(P(v)) = \mathrm{enc_o}(v)$ --- no such law is claimed or wanted; the two
encoders address different structures (a deixis codec must write keys; ontos'
elides them), and the specification records that their byte strings must differ.
ontos bytes are preserved \emph{through the retraction into the ontos model}, which is
the exact guarantee migration needs and the only one on offer.
\end{corollary}

Two remarks make the correspondence architectural rather than curiosities. First,
sibling order: ontos holds order to be part of a tuple, deixis holds order not to be
part of a node --- and Thm.~\ref{thm:kappa}(iii) is why there is no conflict: when a
key spells a position, the position \emph{is} the key, so the ordering lives in key
space and canonical order restates it. Second, recognition never repairs: sparse,
shifted, or non-canonically spelled domains are refused, not normalized ---
by the same identity-coarsening argument as in \S\ref{sec:pos}. The \emph{earlier}
projection into $L(\Bytes)$, its recognizer and laws (i)--(iv) were implemented and \emph{conformance-tested} ---
hand-authored vectors replayed by independent implementations, including replay of
the frozen ontos codec corpus through $R \circ P$; no proof-assistant verification is
claimed. The composition with $E$ above is a mathematical transport; this release
does not claim that the external bridge has adopted the mandatory-value API.\footnote{\texttt{projection/deixis} in the ontos repository; the deixis
floor, the positional spelling, and the set profile are conformance-tested across
Rust, Go, TypeScript, and Python against shared hand-authored vectors --- see
\emph{Artifact availability} below. A machine-checked formalization is future work.}

\section{Related work}\label{sec:related}

Initial-algebra semantics of data types goes back to \cite{gtw}; W-types are
Martin-L\"of's \cite{martinlof}, and their presentation via containers/polynomial
functors is \cite{aag}. Setoids as carriers-with-equality are the standard
constructive treatment of quotients \cite{hofmann}. The mathematical ingredients here
are individually standard --- structural lifting of an equivalence, a recursive path
action, injective framed encodings, an elementary countability characterization ---
and the order-preserving integer trick is folklore in ordered key--value stores;
the contribution claimed is the synthesis and its discipline, not the ingredients.

The closest engineered precedents deserve a direct comparison. Deterministic CBOR
\cite{rfc8949} already specifies self-delimiting values and bytewise ordering of
deterministically encoded map keys; IPLD \cite{ipldcodec} separates a data model
(with maps \emph{and} lists as kinds) from named codecs, and IPLD Schemas
\cite{ipldschema} map richer kinds (structs, unions, enums) onto that model through
representation strategies, with Advanced Data Layouts presenting lensed views over
generic traversal. deixis differs on four axes. \emph{Identity}: deterministic CBOR provides
protocol-selected canonical serialization rules \cite{rfc8949}, and IPLD's abstract
model leaves canonicalization to specific codecs (DAG-CBOR's canonical form and
rejection requirements, with documented decoder leniencies \cite{dagcbor}) --- in
both, identity is whatever byte equality of a chosen encoding induces; deixis first
parameterizes semantic value identity and proves the lift, so equality is defined (and
proved a congruence) before any encoding exists, and canonical bytes are a priced
tier above it. \emph{Container inventory}: IPLD's model carries two containers as
kinds; deixis ships one and recovers the second (and sets, and maps-with-order) as
named profiles with retraction laws --- so ``is this a list?'' is a checkable
judgment with a refusal case, not a kind tag. \emph{Recognition}: representation
strategies specify decodings without algebraic laws attached; profiles carry
retraction, exactness, and identity-transport laws, with hand-authored vectors
pinning the refusals. \emph{Migration}: the ontos bridge states and tests retraction
laws against an incumbent's frozen bytes --- we claim only what is established here,
a proved and vector-tested retraction against one incumbent, and make no assertion
about compatibility results neighboring systems may state for theirs. Content
addressing over trees and DAGs is Merkle's \cite{merkle}; git
and IPLD exploit it with fixed vocabularies, where deixis deliberately ships one
container and receives the rest as embeddings.

The profile interface itself has close formal relatives, and the comparison belongs
at the level of laws. A profile's $(P, R)$ with retraction and exactness is a partial
isomorphism in the sense used by invertible syntax descriptions
\cite{rendel}: one total direction, one partial, composing to identity on the
recognized image. Ornaments \cite{ornaments} likewise present specialized datatypes
as structured views over a common carrier rather than as new types --- the same
motive as profiles-over-one-floor, pursued through indexed descriptions where deixis
uses key spellings. Lenses and prisms \cite{lenses} organize the same territory
around well-behavedness laws; a profile is prism-shaped (build total, match partial),
with exact-or-refusal in place of best-effort match and identity transport as an
added, load-bearing law (Remark~\ref{rem:profilelaws} makes the restriction-idempotent
form precise). The $(P,R)$-with-laws pattern is also lawful \emph{views}
\cite{wadler,mckinnaviews}: the answer to ``are profiles just views?'' is --- views,
plus transport and refusal laws, plus vectors. The Gentle Art of Levitation
\cite{levitation} is the existence proof for Remark~\ref{rem:signal}'s descriptor
language living in the substrate it describes, and its hard-won lesson orders the
roadmap: the descriptor's own identity and decoding must be pinned --- the frozen
codec --- before any self-describing wrapper is sound. For the set profile's
effectivity demands, the Kuratowski-versus-Bishop finiteness distinction
\cite{hottfinsets} is exactly why decidable membership and canonical enumeration are
required; the listable-and-decidable case is the codec-tier case. IPLD's Advanced
Data Layouts present synthesized views over substrate data \cite{ipldadl} and are
the nearest engineered relative; profiles differ in demanding the round-trip and
transport laws as admission criteria, vectors included, rather than as conventions
of a codec implementation.

\paragraph{Artifact availability.}
Release \texttt{v0.1.0} (commit \texttt{7a1193c}) shipped four implementations of
the previous leaf/structure core, the positional spelling and the slot-member set
profile, with 103 black-box checks per implementation. That is historical evidence
for that model. Release \texttt{v0.2.0} supplies the mandatory-value interfaces in
Rust, Go, TypeScript and Python. Its source tag includes this manuscript and the
independently authored \texttt{mnode-*.json} expectations for the mandatory model,
alongside the earlier corpus with its explicit optional-slot interpretation.
The evidence ledger, \texttt{docs/EVIDENCE.md}, records measured counts and tested
revisions. It records codec and downstream adoption status separately.

The harness \texttt{tools/conformance/harness.mjs} drives each implementation as a
subprocess over NDJSON, sending every case \emph{minus its expected judgment}.
No implementation reads the oracle. Continuous integration runs that harness,
per-language format, lint, type and unit checks, specification checks, and two
compilation passes of this manuscript. Those runs are conformance evidence, not
proof. The mathematical derived operations are exercised through the core
accessors; they are not all public methods. The external ontos bridge pins the
v0.2.0 core and position packages in all three of its languages. Its tests run in
ontos's continuous integration; this repository records that adoption without
re-running them.

\section{Conclusion}

deixis fixes the structural floor --- one closed grammar, one open slot, keys as
uninterpreted bytes --- and this paper has stated and proved the core properties the
layers above lean on: lifted identity is a congruence characterized exactly by
pointing; resolution is a lawful partial action; a node is exactly its parts, and
attachment under a fresh key is exact and available at every node; complete-invariant encoders exist
precisely up to countability, computably only strictly less often, and lift through
the combinator; positions embed with order preserved by spelling alone; sequence and
set readings enter as profiles with retraction laws, and the general reading
principle is stated at its proved scope (Remark~\ref{rem:scope}); and the incumbent
value floor above is recovered as one dense instance, exactly, its frozen bytes
preserved through the retraction. What remains open is recorded rather than implied:
a frozen codec with a normative decoder, the scoped profile-universality theorem, and
mechanized proofs. The floor's refusals --- no node kinds, no key meanings, no
repair, structural identity fixed with only its value clause supplied --- are not
austerity for its own sake; each is what makes an eternity of the guarantees above
affordable. \emph{A node is what it points at} is, in the end, both the definition
and the theorem.

\begin{thebibliography}{9}\small\raggedright

\bibitem{gtw} J.~Goguen, J.~Thatcher, E.~Wagner.
\emph{An initial algebra approach to the specification, correctness and
implementation of abstract data types}. In: Current Trends in Programming
Methodology IV, 1978.

\bibitem{martinlof} P.~Martin-L\"of.
\emph{Intuitionistic Type Theory}. Bibliopolis, 1984.

\bibitem{aag} M.~Abbott, T.~Altenkirch, N.~Ghani.
\emph{Containers: constructing strictly positive types}.
Theoretical Computer Science 342(1), 2005.

\bibitem{hofmann} M.~Hofmann.
\emph{Extensional Constructs in Intensional Type Theory}. Springer, 1997.

\bibitem{soare} R.~I.~Soare.
\emph{Recursively Enumerable Sets and Degrees}. Springer, 1987.

\bibitem{cockettlack} J.~R.~B.~Cockett, S.~Lack.
\emph{Restriction categories I: categories of partial maps}.
Theoretical Computer Science 270(1--2), 2002.

\bibitem{wadler} P.~Wadler.
\emph{Views: a way for pattern matching to cohabit with data abstraction}.
POPL 1987.

\bibitem{mckinnaviews} C.~McBride, J.~McKinna.
\emph{The view from the left}. J.\ Functional Programming 14(1), 2004.

\bibitem{levitation} J.~Chapman, P.-E.~Dagand, C.~McBride, P.~Morris.
\emph{The gentle art of levitation}. ICFP 2010.

\bibitem{aagm} M.~Abbott, T.~Altenkirch, N.~Ghani, C.~McBride.
\emph{Constructing polymorphic programs with quotient types}. MPC 2004.

\bibitem{hottfinsets} D.~Frumin, H.~Geuvers, L.~Gondelman, N.~van der Weide.
\emph{Finite sets in homotopy type theory}. CPP 2018.

\bibitem{elias} P.~Elias.
\emph{Universal codeword sets and representations of the integers}.
IEEE Trans.\ Information Theory 21(2), 1975.

\bibitem{rendel} T.~Rendel, K.~Ostermann.
\emph{Invertible syntax descriptions: unifying parsing and pretty printing}.
Haskell Symposium, 2010.

\bibitem{ornaments} P.-E.~Dagand, C.~McBride.
\emph{A categorical treatment of ornaments}. LICS 2013; arXiv:1212.3806.

\bibitem{lenses} J.~N.~Foster, M.~B.~Greenwald, J.~T.~Moore, B.~C.~Pierce,
A.~Schmitt. \emph{Combinators for bidirectional tree transformations}.
ACM TOPLAS 29(3), 2007.

\bibitem{rfc8949} C.~Bormann, P.~Hoffman.
\emph{Concise Binary Object Representation (CBOR)}. RFC 8949, IETF, December 2020.

\bibitem{merkle} R.~Merkle.
\emph{A digital signature based on a conventional encryption function}.
CRYPTO 1987.

\bibitem{ipldcodec} Protocol Labs.
\emph{IPLD documentation: Codecs}. \texttt{https://ipld.io/docs/codecs/}
(accessed August 2026).

\bibitem{ipldschema} Protocol Labs.
\emph{IPLD Schemas: representation strategies}.
\texttt{https://ipld.io/docs/schemas/features/representation-strategies/}
(accessed August 2026).

\bibitem{ipldadl} Protocol Labs.
\emph{IPLD: Advanced Data Layouts}.
\texttt{https://ipld.io/docs/advanced-data-layouts/} (accessed August 2026).

\bibitem{dagcbor} Protocol Labs.
\emph{DAG-CBOR codec specification}.
\url{https://ipld.io/specs/codecs/dag-cbor/spec/} (accessed August 2026).

\end{thebibliography}

\end{document}
